\section*{Chapter \ref{ch:g10:specifications-and-tests} --- Specifications, Assertions and Tests}

\begin{solution}{exo:g10:specifications-and-tests:1}
Both examples fail. \code{doctest} compares the \emph{text} that the call prints with the expected text:
\code{half(10)} prints \code{5.0} (the operator \code{/} always gives a decimal number), not \code{5}; \code{half(7)}
prints \code{3.5}, not \code{3}. Either the examples are wrong (the expected answers should be \code{5.0} and
\code{3.5}) or the body is (\code{n // 2} gives \code{5} and \code{3}); the \omterm{def:g10:specifications-and-tests:specification}{specification} must decide which.
\end{solution}

\begin{solution}{exo:g10:specifications-and-tests:2}
\code{3 > 5} is false, so \code{loud(1)} is not evaluated: the line prints \code{False}. \code{3 < 5} is true, so
\code{loud(2)} is not evaluated: \code{True}. For \code{loud(0) or loud(4)}: \code{loud(0)} prints
\code{evaluating loud(0)} and returns \code{0}, which counts as false, so \code{loud(4)} is evaluated too and prints
\code{evaluating loud(4)}; the result of \code{or} is its last evaluated operand, so the line prints \code{4}.
\end{solution}

\begin{solution}{exo:g10:specifications-and-tests:3}
``Return the smallest strictly positive number of \code{t}. \omterm{def:g10:specifications-and-tests:pre-post}{Precondition}: \code{t} is a list of numbers that
contains at least one number greater than $0$. \omterm{def:g10:specifications-and-tests:pre-post}{Postcondition}: the result is an element of \code{t}, it is greater
than $0$, and no element of \code{t} that is greater than $0$ is smaller than it.'' Asking only for a non-empty list
is not enough: \code{[-3, 0]} is not empty and has no positive element.
\end{solution}

\begin{solution}{exo:g10:specifications-and-tests:4}
Three kinds: \code{x < lo}, \code{lo <= x <= hi}, \code{x > hi}; boundaries \code{x = lo} and \code{x = hi}; and the
edge of the \omterm{def:g10:specifications-and-tests:pre-post}{precondition} \code{lo = hi}. Tests (with \code{lo = 0}, \code{hi = 10} unless stated):
\code{clamp(-5, 0, 10)} gives \code{0}; \code{clamp(0, 0, 10)} gives \code{0}; \code{clamp(4, 0, 10)} gives \code{4};
\code{clamp(10, 0, 10)} gives \code{10}; \code{clamp(12, 0, 10)} gives \code{10}; \code{clamp(7, 3, 3)} gives
\code{3}.
\end{solution}

\begin{solution}{exo:g10:specifications-and-tests:5}
The \omterm{def:g10:specifications-and-tests:test-case}{boundary case} \code{is_adult(18)}: the expected result is \code{True}, the function returns \code{False}. The
body must be \code{return age >= 18} (item 5 of the checklist, the wrong inequality).
\end{solution}

\begin{solution}{exo:g10:specifications-and-tests:6}
Removing an element while the loop walks through the same list shifts the elements that follow it: after
\code{-1} at position $0$ is removed, \code{-2} moves to position $0$ and the loop goes on at position $1$, so
\code{-2} is never examined. The two bugs: the loop skips elements, and the function changes the caller's list (a
\omterm{def:g10:specifications-and-tests:side-effect}{side effect} the \omterm{def:g10:specifications-and-tests:specification}{specification} does not announce). The fix, in \code{exercises/remove_negatives.py}:
\code{result = []}, then \code{for x in t:} \code{if x >= 0: result.append(x)}, and \code{return result}. On
\code{[-1, -2, 3]} it returns \code{[3]} and leaves the argument unchanged (checked by the chapter's tests).
\end{solution}

\begin{solution}{exo:g10:specifications-and-tests:7}
Kinds: the seven months of $31$ days, the four of $30$, and February, which splits into the four kinds of year.
Boundaries: month \code{1} and month \code{12}. The suite in \code{exercises/days_in_month.py} has seven tests:
January and December $2025$ ($31$), April $2025$ ($30$), February in $2025$ ($28$), $2024$ ($29$), $1900$ ($28$) and
$2000$ ($29$). The body: \code{if month == 2: return 29 if leap_year(year) else 28}; then
\code{if month in (4, 6, 9, 11): return 30}; else \code{return 31}. The chapter's tests compare it with Python's
\code{calendar} module on every month of fifty-four years.
\end{solution}

\begin{solution}{exo:g10:specifications-and-tests:8}
\code{[]} passes (\code{True}), \code{[4]} passes (\code{True}), \code{[1, 2, 3]} passes (\code{True}),
\code{[3, 2, 1]} passes (\code{False}, as expected), \code{[1, 5, 2, 6]} fails: the function returns \code{True}
because $1 \le 6$, but the list is not sorted. The author forgot lists that go down somewhere \emph{inside} while
their ends are in order; the fix compares every pair of neighbours.
\end{solution}

\begin{solution}{exo:g10:specifications-and-tests:9}
Take a function \code{bad_is_prime(n)} that returns \code{False} when \code{n == 1009} and
\code{is_prime(n)} otherwise. It passes every
test from $2$ to $1000$ ($1009$ is not among them) and is wrong, since $1009$ is prime. That is the proof of the
proposition on an example: passing tests says nothing about the inputs that were not tried. To know that
\code{is_prime} is correct for every \code{n}, one needs a proof that its body implements the definition of a prime,
the kind of argument studied in grade 11.
\end{solution}

\begin{solution}{exo:g10:specifications-and-tests:10}
\code{sign(0)} reaches the end of the function without a \code{return} and returns \code{None}: item 4, the
non-exhaustive \code{if}. Adding \code{else: return 0} (\code{exercises/sign.py}) makes the three branches cover every
number.
\end{solution}

\begin{solution}{exo:g10:specifications-and-tests:11}
In \code{exercises/percent.py}, three \omterm{def:g10:specifications-and-tests:assertion}{assertions} surround the computation
\code{result = 100 * part / whole}: before it, \code{assert whole > 0} with the message ``whole must be
positive'' and \code{assert 0 <= part <= whole} with the message ``part must lie between 0 and whole''; after it,
\code{assert 0 <= result <= 100}; then \code{return result}. The tests check \code{percent(1, 4) == 25} and that
\code{(1, 0)}, \code{(-1, 4)} and \code{(5, 4)} each raise \code{AssertionError}.
\end{solution}

\begin{solution}{exo:g10:specifications-and-tests:12}
\emph{(a)} Each \code{if} is taken or not, independently: $2 \times 2 = 4$ paths. \emph{(b)} One test makes every
line run: a child on a Tuesday ($10 \to 5 \to 3$). \emph{(c)} Four tests follow every path: an adult on a Monday
($10$), an adult on a Tuesday ($8$), a child on a Monday ($5$), a child on a Tuesday ($3$). Running every line is
weaker than following every path: the single test of (b) would not notice a bug that only shows for an adult on a
Tuesday.
\end{solution}

\begin{solution}{exo:g10:specifications-and-tests:13}
In \code{exercises/roman.py}: a table of the thirteen values \code{1000 M}, \code{900 CM}, \code{500 D},
\code{400 CD}, \code{100 C}, \code{90 XC}, \code{50 L}, \code{40 XL}, \code{10 X}, \code{9 IX}, \code{5 V},
\code{4 IV}, \code{1 I}, from the largest down; for each value, while \code{n >= value}, append its letters and
subtract it. The \omterm{def:g10:specifications-and-tests:pre-post}{precondition} is checked by \code{assert 1 <= n <= 3999}. The suite tests \code{roman(1)},
\code{4}, \code{9}, \code{40}, \code{900}, \code{2026} (\code{MMXXVI}), \code{3999} (\code{MMMCMXCIX}), and that
\code{0} and \code{4000} raise \code{AssertionError}; the chapter's tests also check that the $3999$ numerals are all
different.
\end{solution}

\begin{solution}{exo:g10:specifications-and-tests:14}
Item 6, equality between decimal numbers: $0.1$ is not stored exactly, so after ten additions \code{x} is very
close to $1$ but not equal to \code{1.0}, and the condition \code{x != 1.0} stays true for ever. First fix:
\code{for steps in range(1, 11): x = x + 0.1}, a loop that counts its steps instead of testing \code{x}. Second fix
(\code{exercises/count_up.py}): \code{while not math.isclose(x, 1.0)}, which stops after $10$ steps (checked by the
chapter's tests).
\end{solution}

\begin{solution}{exo:g10:specifications-and-tests:15}
\emph{(a)} survives: with \code{100} in the first test, $2000$ still returns \code{True} there and $2024$ is not a
century. \emph{(b)} survives: neither $2024$ nor $2000$ reaches the line \code{return False}. \emph{(c)} is killed:
$2024$ then gives \texttt{2024 \% 400 == 0}, false. \emph{(d)} is killed: $2024$ gives \code{24 != 0}, true, so the
mutant returns \code{False}. Adding $1900$ kills both survivors: mutant (a) returns \code{True} for it at the first
test, mutant (b) returns \code{True} at the second.
\end{solution}

\begin{solution}{pb:g10:specifications-and-tests:1}
\textbf{1.} \code{year} is a whole number, $1583$ or later: the rule belongs to a calendar instituted in $1582$, and
asking about earlier years makes no sense for it.

\textbf{2.} Not divisible by $4$: not leap. Divisible by $4$ but not by $100$: leap. Divisible by $100$ but not by
$400$: not leap. Divisible by $400$: leap.

\textbf{3.} The century years are $1600, 1700, \dots, 2400$: $9$ of them. The leap ones are $1600$, $2000$ and
$2400$: $3$.

\textbf{4.} Sam's version is wrong exactly on the $9 - 3 = 6$ centuries that are not leap. From $1583$ to $2400$
there are $2400 - 1583 + 1 = 818$ years, so the fraction is $6/818$, about $0.7\,\%$.

\textbf{5.} The suite of four \omterm{def:g10:specifications-and-tests:test-case}{test cases}: $2023$ gives \code{False}, $2024$ gives \code{True}, $1900$ gives
\code{False}, $2000$ gives \code{True}.

\textbf{6.} Each year belongs to exactly one of the four kinds, so three years can cover at most three kinds.

\textbf{7.} It fails the test of $1900$. To let the bug through, a suite must contain no year of the third kind
(divisible by $100$ and not by $400$).

\textbf{8.} $1583$ is the smallest year the \omterm{def:g10:specifications-and-tests:pre-post}{precondition} allows, the edge of the allowed inputs; it is not divisible
by $4$, so the expected result is \code{False}.

\textbf{9.} Three \code{==} (lines $8$, $10$, $12$): $3$; one \code{True} and one \code{False}: $2$; \code{400},
\code{100} and \code{4} have two replacements each: $6$; three \code{0} on the right of \code{==}: $3$. In all
$3 + 2 + 6 + 3 = 14$ mutants.

\textbf{10.} The mutant reads \texttt{if year \% 100 != 0: return False}. On $2024$ the condition is true
(\texttt{2024 \% 100} is $24$), so it returns \code{False} instead of \code{True}; the test of $2024$ fails.

\textbf{11.} The first line becomes \texttt{if year \% 4 == 0: return True}: every year divisible by $4$ is leap, so it
gets the third kind wrong, exactly like Sam's version ($1900$ gives \code{True}).

\textbf{12.} The program prints \code{boundary suite (4 years) kills 14}: all fourteen.

\textbf{13.} The mutant returns \code{False} at the first test for every multiple of $400$; the suite contains
$2000$, whose expected result is \code{True}, so that test fails.

\textbf{14.} $9$ century years among $818$ is a frequency of about $9/818 \approx 0.011$, so about one century among
$100$ random years (and quite possibly none).

\textbf{15.} The survivors are: line 9 \code{True} to \code{False}; line 11 \code{False} to \code{True}; line 8
\code{400} to \code{100} and \code{400} to \code{4}; line 10 \code{100} to \code{400}; the \code{0} of line 8 and of
line 10 changed to \code{1}. All seven change only what happens on the first two tests, the ones about centuries;
every year that is not a century reaches the last line, which is intact in all seven.

\textbf{16.} None: the seeded random suite contains no year divisible by $100$, so the lines about centuries never
decide its results, and a change to them is invisible.

\textbf{17.} A century that is not leap ($1900$, say) kills five survivors: those of line 11, the two changes of
\code{400}, the change of \code{100} and the \code{0} of line 10; each makes $1900$ leap. A multiple of $400$
($2000$) kills the other two (line 9, and the \code{0} of line 8). No single year kills all seven: the first group
needs a year divisible by $100$ and not by $400$, the second a year divisible by $400$.

\textbf{18.} By question 17, two years are needed and two suffice: $1900$ and $2000$.

\textbf{19.} The four years chosen from the \omterm{def:g10:specifications-and-tests:specification}{specification} kill all $14$ mutants; the $100$ random years kill only $7$.
Four well-chosen tests beat a hundred random ones, because the random ones never land on the rare kinds of input
where the \omterm{def:g10:specifications-and-tests:specification}{specification} changes its answer.
\end{solution}
